\begin{proposition}[4.6.15]
\label{IV.2.4.6.15-fr}
\begin{enumerate}
\item[\textup{(i)}] Sean $k$ un cuerpo, $X$ un $k$-preesquema de tipo finito. Para que
$X$ sea geométricamente puntualmente íntegro, es necesario y suficiente que
sea geométricamente reducido y que el número geométrico de sus componentes
irreducibles sea igual al número geométrico de sus componentes conexas.
\item[\textup{(ii)}] Sea $X$ un preesquema localmente de tipo finito sobre $k$. Si $X$
es geométricamente puntualmente íntegro en un punto $x$, existe una vecindad
abierta $U$ de $x$ que es geométricamente puntualmente íntegra. En otras
palabras, el conjunto de los $x\in X$ en los que $X$ es geométricamente
puntualmente íntegro es abierto en $X$.
\end{enumerate}
\end{proposition}

\begin{proof}
\begin{enumerate}
\item[(i)] Dada una extensión algebraicamente cerrada $\Omega$ de $k$,
$X_{(\Omega)}$ es puntualmente íntegro (o localmente íntegro, lo que equivale a
lo mismo puesto que $X_{(\Omega)}$ es noetheriano) si y solo si es reducido y
si el número de sus componentes conexas es igual al número de sus componentes
irreducibles (\textbf I, 6.1.10). Basta entonces aplicar (4.5.1), (4.6.1) y la primera
observación de (4.6.14).

\item[(ii)] Puesto que la cuestión es local sobre $X$, es posible limitarse al
caso en que $X$ es un preesquema de tipo finito sobre $k$. Además, se sabe que
existe una vecindad abierta de $x$ en $X$ que es geométricamente reducida
(4.6.13), por lo que puede suponerse que $X$ es geométricamente reducido.
Existe entonces una extensión finita $k'$ de $k$ tal que
$X'=X\otimes_k k'$ sea reducido y que sus componentes irreducibles sean
geométricamente irreducibles (4.6.8). Sea $p:X'\to X$ la proyección canónica,
que es un morfismo finito sobreyectivo, y sean $x'_j$
($1\leqslant j\leqslant r$) los puntos de $p^{-1}(x)$~; por hipótesis, $X'$ es
íntegro en cada uno de los $x'_j$, por lo que cada $x'_j$ tiene una vecindad
abierta $V'_j$ en $X'$ que es íntegra. Como $p$ es cerrado
(\textbf{II}, 6.1.10), cada
uno de los $p(V'_j)$ es una vecindad de $x$, por lo que existe una vecindad abierta
$U$ de $x$ tal que $p^{-1}(U)$ está contenido en la unión de los $V'_j$, y es,
por consiguiente, íntegro en cada uno de sus puntos, y por tanto localmente
íntegro. Las componentes irreducibles de $p^{-1}(U)$ son, pues, abiertas y
disjuntas dos a dos, y como son geométricamente irreducibles (4.5.9),
$p^{-1}(U)=U\otimes_k k'$ es geométricamente puntualmente íntegro en virtud de
(i), y lo mismo ocurre, por tanto, con $U$ por definición.
\end{enumerate}
\end{proof}

\begin{proposition}[4.6.16]
\label{IV.2.4.6.16-fr}
Sean $k$ un cuerpo, $X$ un $k$-preesquema localmente noetheriano, $\mathscr F$
un $\mathcal O_X$-módulo coherente. Las condiciones siguientes son
equivalentes~:
\begin{enumerate}
\item[\textup{a)}] Para toda extensión de tipo finito $k'$ de $k$, si se pone
$X'=X\otimes_k k'$, entonces $\mathscr F'=\mathscr F\otimes_k k'$ es un
$\mathcal O_{X'}$-módulo reducido \textup{(3.2.2)}.
\item[\textup{b)}] Para toda extensión finita radicial $k'$ de $k$, $\mathscr F'$ es un
$\mathcal O_{X'}$-módulo reducido.
\item[\textup{c)}] $\mathscr F$ es reducido, y si $\mathscr J$ es el ideal de
$\mathcal O_X$ anulador de $\mathscr F$, el subpreesquema cerrado definido por
$\mathscr J$ es geométricamente reducido sobre $k$.
\end{enumerate}

\oldpage[IV-2]{73}

Además, si $X$ es localmente de tipo finito sobre $k$, estas condiciones son
también equivalentes a la siguiente~:
\begin{enumerate}
\item[\textup{d)}] Para toda extensión $k'$ de $k$ \textup{(o} para \emph{una} extensión
algebraicamente cerrada $k'$ de $k$), $\mathscr F'$ es un
$\mathcal O_{X'}$-módulo reducido.
\end{enumerate}
\end{proposition}

\begin{proof}
Es claro que \emph{a)} implica \emph{b)}. La condición \emph{b)} implica ante todo que $\mathscr F$
es reducido~: esto significa (3.2.3) que, si $Y$ es el subpreesquema cerrado de
$X$ definido por el ideal coherente $\mathscr J$ de $\mathcal O_X$, anulador de
$\mathscr F$, $Y$ es \emph{reducido}, y existe un $\mathcal O_Y$-módulo coherente sin
torsión $\mathscr G$ de rango~1 sobre toda componente irreducible de $Y$, tal
que $j_*(\mathscr G)=\mathscr F$, siendo $j:Y\to X$ la inyección canónica.
Ahora bien, para \emph{toda} extensión $k'$ de $k$, el anulador de $\mathscr F'$ es
$\mathscr J'=\mathscr J\otimes_k k'$, pues el morfismo de proyección
$X'\to X$ es plano (2.1.11)~; el subpreesquema cerrado de $X'$ definido por
$\mathscr J'$ es $Y'=Y\otimes_k k'$, y si $j':Y'\to X'$ es la inyección
canónica y $\mathscr G'=\mathscr G\otimes_k k'$, se tiene
$\mathscr F'=j'_*(\mathscr G')$. Observemos también que, para toda extensión
$k'$ de $k$ tal que $X'$ sea localmente noetheriano, $\mathscr F'$ carece de
ciclos primos asociados inmersos (4.2.7). Además, todo punto maximal $y'$ de
$Y'$ está sobre un punto maximal $y$ de $Y$~; como $\mathscr G_y$ es isomorfo
a $\mathcal O_{Y,y}$, $\mathscr G'_{y'}$ es isomorfo a
$\mathcal O_{Y',y'}$. Se ve, por tanto (por (3.2.3)), que equivale a decir que
$\mathscr F'$ es reducido o que el preesquema $Y'$ es reducido, cuando $X'$ es
\emph{localmente noetheriano}. El hecho de que \emph{b)} implique \emph{c)} y de que \emph{c)} implique \emph{a)}
es, por tanto, consecuencia de (4.6.1, \emph{d)} y \emph{b)})~; en \emph{a)}, la hipótesis de que
$k'$ es una extensión de tipo finito de $k$ solo se formula para asegurar que
$X'$ sea localmente noetheriano. Cuando $X$ es localmente de tipo finito,
$X'$ es localmente noetheriano para toda extensión $k'$ de $k$, lo que
demuestra en este caso la equivalencia de \emph{d)} y las demás condiciones.
\end{proof}

\begin{definition}[4.6.17]
\label{IV.2.4.6.17-fr}
Cuando $\mathscr F$ satisface las condiciones equivalentes de \textup{(4.6.16)}, se
dice que $\mathscr F$ es geométricamente reducido sobre $k$, o
separable sobre $k$.
\end{definition}

\begin{proposition}[4.6.18]
\label{IV.2.4.6.18-fr}
Sean $k$ un cuerpo, $X$ un $k$-preesquema localmente noetheriano, $\mathscr F$
un $\mathcal O_X$-módulo coherente. Las condiciones siguientes son
equivalentes~:
\begin{enumerate}
\item[\textup{a)}] Para toda extensión de tipo finito $k'$ de $k$, si se pone
$X'=X\otimes_k k'$, entonces $\mathscr F'=\mathscr F\otimes_k k'$ es un
$\mathcal O_{X'}$-módulo íntegro \textup{(3.2.4)}.
\item[\textup{b)}] Para toda extensión finita $k'$ de $k$, $\mathscr F'$ es un
$\mathcal O_{X'}$-módulo íntegro.
\item[\textup{c)}] $\mathscr F$ es reducido (o íntegro), y si $\mathscr J$ es el ideal
de $\mathcal O_X$ anulador de $\mathscr F$, el subpreesquema cerrado definido
por $\mathscr J$ es geométricamente íntegro.
\end{enumerate}
Además, si $X$ es localmente de tipo finito sobre $k$, estas condiciones son
también equivalentes a la siguiente~:
\begin{enumerate}
\item[\textup{d)}] Para toda extensión $k'$ de $k$ \textup{(o} para \emph{una} extensión
algebraicamente cerrada $k'$ de $k$), $\mathscr F'$ es un
$\mathcal O_{X'}$-módulo íntegro.
\end{enumerate}
\end{proposition}

\begin{proof}
Las condiciones \emph{a)}, \emph{b)} y \emph{c)} implican que $\mathscr F$ es íntegro, por lo que
$\operatorname{Ass}(\mathscr F)$ se reduce a un solo punto $x$, que es el punto
genérico de $\operatorname{Supp}(\mathscr F)$. Para \emph{toda} extensión $k'$ de $k$
para la cual $X'$ sea localmente noetheriano, los ciclos primos asociados a
$\mathscr F'$ son, por tanto, no inmersos, y son los puntos maximales de
$\operatorname{Supp}(\mathscr F')$, todos ellos situados sobre $x$. Por otra
parte, en virtud de (4.6.16), cada una de las condiciones \emph{a)}, \emph{b)}, \emph{c)} implica que
$\mathscr F$ es geométricamente reducido~; es trivial que \emph{a)} implica \emph{b)}, y \emph{b)}
implica que $\operatorname{Supp}(\mathscr F)$ es geométricamente irreducible
(4.5.9), y por tanto \emph{c)}, puesto que el subpreesquema $Y$ definido

\oldpage[IV-2]{74}

por $\mathscr J$ es entonces geométricamente reducido y geométricamente
irreducible, y por tanto geométricamente íntegro. Recíprocamente, si $Y$ es
geométricamente íntegro, para \emph{toda} extensión $k'$ de $k$,
$\operatorname{Supp}(\mathscr F')$ es irreducible, por lo que, si $X'$ es
localmente noetheriano, $\mathscr F'$ es íntegro~; esto muestra que \emph{c)} implica
\emph{a)}, y que \emph{c)} implica \emph{d)} cuando $X$ es localmente de tipo finito sobre $k$.
\end{proof}
